%% cover_letter.tex — C9: SVD Interlacing / Physics-Reduction Proposition
%% SIAM Journal on Matrix Analysis and Applications
\documentclass[11pt,a4paper]{letter}
\usepackage[T1]{fontenc}
\usepackage[utf8]{inputenc}
\usepackage[margin=2.5cm]{geometry}
\usepackage{microtype}
\usepackage{amsmath}

\signature{Anoop~V.~Eluvathingal\\
           Ph.D.\ Candidate, Dept.\ of Electrical Engineering\\
           Indian Institute of Technology Madras\\
           Chennai 600036, India\\[2pt]
           \texttt{arun.eluvathingal@iitm.ac.in}}

\address{Department of Electrical Engineering\\
         Indian Institute of Technology Madras\\
         Chennai 600036, India}

\begin{document}
\begin{letter}{Editor-in-Chief\\
               \textit{SIAM Journal on Matrix Analysis and Applications}}

\opening{Dear Editor,}

We submit for your consideration the manuscript:

\begin{center}
  \textbf{Singular Value Interlacing for Physics-Structured Projections:
  Condition-Number Inheritance and the Physics-Reduction Proposition}
\end{center}

\textbf{Mathematical content.}
This paper proves three results in matrix analysis.

\emph{Theorem~1} establishes Cauchy-type SVD interlacing inequalities
for a physics-structured projection $\tilde{\mathbf{H}} = \mathbf{H}\mathbf{P}$,
where $\mathbf{P}$ is determined by the algebraic constraints of a
differential-algebraic equation (DAE) system:
\[
  \sigma_k(\mathbf{H}) \geq \tilde{\sigma}_k \geq
  \frac{\sigma_{k+n_a}(\mathbf{H})}{\sqrt{1+\rho_Q^2}},
  \quad k = 1,\ldots,r.
\]
The correction factor $1/\sqrt{1+\|\mathbf{Q}\|^2}$ is new; standard
Cauchy interlacing (for column-restriction by an orthonormal matrix)
does not carry this term, but the physics projection $\mathbf{P}$
generally does not have orthonormal columns.

\emph{Theorem~2} derives an explicit condition-number inheritance bound:
\[
  \kappa(\tilde{\mathbf{H}}) \leq
  \kappa(\mathbf{H}) \cdot \frac{1+\rho}{1-\rho} \cdot \Delta_r^{-1},
\]
where $\rho$ is the algebraic coupling ratio and $\Delta_r$ is the
spectral gap.  When $\Delta_r \geq (1+\rho)/(1-\rho)$, the condition
number is inherited exactly: $\kappa(\tilde{\mathbf{H}}) \leq \kappa(\mathbf{H})$.

\emph{Proposition~1 (Physics-Reduction Proposition)} specialises
Theorem~2 to the 6-state IBR fault-current Jacobian with $r=2$
dynamic states, proving that $\kappa(\mathbf{H}) < 30$ implies
$\kappa(\tilde{\mathbf{H}}) < 30$ for standard generator parameters.

All bounds are shown to be \emph{sharp}: we exhibit explicit matrix
families attaining each inequality.

\textbf{Novelty.}
Prior work on SVD interlacing \cite{Thompson1972} addresses
column-submatrix restrictions.  The physics-structured projection
$\mathbf{P} = (\mathbf{I}; \mathbf{Q})$ arising from DAE constraint
elimination is a perturbation of a column restriction and requires a
separate analysis.  The condition-number inheritance bound and
the spectral-gap criterion (Corollary~1) are new.

\textbf{AMS classification.} 15A18 (eigenvalues), 15A60 (norms),
65F35 (numerical linear algebra), 93B30 (identifiability).

\textbf{Suggested reviewers.}
\begin{itemize}
  \item Prof.\ Roger Horn (University of Utah) --- matrix analysis
  \item Prof.\ Ilse Ipsen (NC State University) --- perturbation theory,
        condition numbers
  \item Prof.\ Zlatko Drmac (University of Zagreb)
        --- SVD algorithms and perturbation bounds
  \item Prof.\ Ran Ran (Beijing Institute of Technology)
        --- interlacing and eigenvalue inequalities
\end{itemize}

This manuscript is not under review elsewhere.

\closing{Sincerely,}
\end{letter}
\end{document}
